\documentclass{cskkuproject}

% ==========================================
% ข้อมูลโครงงาน — กรอกข้อมูลที่นี่
% ==========================================

\titleTH{ชื่อโครงงานภาษาไทย}
\titleEN{Project Title in English}

% ข้อมูลนักศึกษาคนที่ 1
\firstStudentID{6xxxxxxxxx-x}
\firstStudentPrefixTH{นาย}
\firstStudentPrefixEN{Mr.}
\firstStudentNameTH{ชื่อ นามสกุล}
\firstStudentNameEN{Firstname Lastname}

% ข้อมูลนักศึกษาคนที่ 2
% *** ถ้าทำโครงงานคนเดียว ให้เว้นว่างไว้ทั้งหมด (เหมือนด้านล่าง) ***
% *** ถ้ามี 2 คน ให้กรอกข้อมูลเหมือนคนที่ 1 ***
\secondStudentID{}
\secondStudentPrefixTH{}
\secondStudentPrefixEN{}
\secondStudentNameTH{}
\secondStudentNameEN{}

% ข้อมูลหลักสูตร
\degreeTH{วิทยาศาสตรบัณฑิต}
\degreeEN{Bachelor of Science}
\majorTH{วิทยาการคอมพิวเตอร์}
\majorEN{Computer Science}

% อาจารย์ที่ปรึกษา
\advisorTH{ผศ.ดร.ชื่อ นามสกุล}
\advisorEN{Asst.Prof.Firstname Lastname, Ph.D.}

% ปีการศึกษา
\academicYearTH{2568}
\academicYearEN{2025}

% รหัสวิชา (CP354 771 = โครงงาน 1, CP354 772 = โครงงาน 2)
\courseCode{CP354 772}
\courseName{โครงงานคอมพิวเตอร์ 2}
\semester{2}

% คำสำคัญ
\keywordsTH{คำสำคัญ 1, คำสำคัญ 2, คำสำคัญ 3}
\keywordsEN{Keyword 1, Keyword 2, Keyword 3}


\begin{document}

% ==========================================
% หน้าปก
% ==========================================
% คำสั่งด้านล่างสร้างหน้าปกอัตโนมัติจากข้อมูลที่กรอกด้านบน
% ไม่ต้องแก้ไขส่วนนี้
\maketitlepage

% ตั้งค่าเลขหน้าส่วนต้นเป็นอักษรไทย (ก, ข, ค, ...)
\pagenumbering{arabic}
\setcounter{page}{0}
\renewcommand{\thepage}{\ThaiAlph{\value{page}}}

% ==========================================
% บทคัดย่อ
% ==========================================
\begin{abstractTH}
% บทคัดย่อภาษาไทย
% หมายเหตุ: ไม่ควรเกิน 300 คำ
% ควรครอบคลุม: วัตถุประสงค์ → วิธีการ → ผลลัพธ์ → การประยุกต์ใช้
%
% *** ลบข้อความตัวอย่างด้านล่างทั้งหมด แล้วเขียนบทคัดย่อของตนเอง ***

โครงงานฉบับนี้มีวัตถุประสงค์เพื่อ\ldots{} โดย\ldots{} ระบบที่พัฒนาขึ้นสามารถ\ldots{} ผลการทดสอบพบว่า\ldots{} ดังนั้น โครงงานนี้จึงเป็นประโยชน์ในด้าน\ldots{}

\end{abstractTH}

\begin{abstractEN}
% ABSTRACT (English)
% Note: Should not exceed 300 words.
% Should cover: Objective, Methodology, Results, Applications
% Use Past tense for objectives and methodology
% Use Present tense for results and applications
%
% *** Delete the example text below and write your own abstract ***

This project aims to\ldots{} The system was developed using\ldots{} The developed system is capable of\ldots{} The experimental results show that\ldots{} Therefore, this project is beneficial for\ldots{}

\end{abstractEN}

% ==========================================
% กิตติกรรมประกาศ
% ==========================================
\begin{acknowledgement}
\input{acknowledgement}
\end{acknowledgement}

% ==========================================
% สารบัญ
% ==========================================
\tableofcontents
\clearpage
\listoffigures
\clearpage
\listoftables
\clearpage

% ==========================================
% สัญลักษณ์และคำย่อ (ถ้ามี — uncomment ด้านล่าง)
% ==========================================
% \begin{symbolpage}
% \input{symbols}
% \end{symbolpage}

% ==========================================
% เนื้อเรื่อง 5 บท
% ==========================================
% เริ่มนับเลขหน้าใหม่เป็นเลขอารบิก (1, 2, 3, ...)
% ไม่ต้องแก้ไขส่วนนี้
\pagenumbering{arabic}
\setcounter{page}{1}

\chapterone
\chaptertwo
\chapterthree
\chapterfour
\chapterfive

% ==========================================
% เอกสารอ้างอิง
% ==========================================
% *** ไม่ต้องแก้ไขส่วนนี้ ***
% ระบบจะดึงรายการอ้างอิงจากไฟล์ references.bib อัตโนมัติ
% ให้แก้ไขรายการอ้างอิงในไฟล์ references.bib แทน
\clearpage
\addcontentsline{toc}{chapter}{เอกสารอ้างอิง}
\makeatletter
\let\ps@plain\ps@empty
\makeatother
\bibliographystyle{ieeetr}
\bibliography{references}

% ==========================================
% ภาคผนวก
% ==========================================
\input{appendices}

% ==========================================
% ประวัติผู้เขียน
% ==========================================
\begin{biography}
% ประวัติผู้เขียน
% หมายเหตุ: ไม่เกิน 1 หน้ากระดาษ
%
% *** ลบข้อความตัวอย่างด้านล่าง แล้วเขียนประวัติของตนเอง ***

\section*{นักศึกษาคนที่ 1}

นาย/นางสาว\ldots{} เกิดเมื่อวันที่\ldots{} ที่\ldots{}
สำเร็จการศึกษาระดับมัธยมศึกษาตอนปลาย จากโรงเรียน\ldots{} เมื่อปีการศึกษา\ldots{}
ปัจจุบันกำลังศึกษาระดับปริญญาตรี หลักสูตรวิทยาศาสตรบัณฑิต
สาขาวิชาวิทยาการคอมพิวเตอร์ วิทยาลัยการคอมพิวเตอร์ มหาวิทยาลัยขอนแก่น

% ถ้ามีรางวัลหรือทุนการศึกษา:
% ได้รับ ... จาก ... เมื่อ พ.ศ. ...

% ถ้ามีประสบการณ์ทางวิชาการ:
% มีประสบการณ์ ...


% --- uncomment ด้านล่างนี้ ถ้ามีนักศึกษาคนที่ 2 ---
% \section*{นักศึกษาคนที่ 2}
%
% นาย/นางสาว\ldots{} เกิดเมื่อวันที่\ldots{} ที่\ldots{}
% สำเร็จการศึกษาระดับมัธยมศึกษาตอนปลาย จากโรงเรียน\ldots{} เมื่อปีการศึกษา\ldots{}
% ปัจจุบันกำลังศึกษาระดับปริญญาตรี หลักสูตรวิทยาศาสตรบัณฑิต
% สาขาวิชาวิทยาการคอมพิวเตอร์ วิทยาลัยการคอมพิวเตอร์ มหาวิทยาลัยขอนแก่น

\end{biography}

\end{document}
